% Methods and Results sections for the autonomous quadrotor landing manuscript.
% IEEEtran class. Requires: amsmath, amssymb, graphicx, booktabs, siunitx.
% Figure/table labels referenced here: fig:architecture, fig:statemachine,
% fig:offsets, tab:params, tab:results.

\section{Methods}

\subsection{System Description}

The landing system estimates the pose of a moving platform from a downward-facing
camera and regulates the vehicle to a touchdown point on that platform using
velocity setpoints. It comprises three elements: an altitude estimator that fuses
two dissimilar range sources, a visual alignment law that acts on image-plane
error, and a supervisory state machine that sequences alignment, descent, and
flare. Fig.~\ref{fig:architecture} shows the arrangement.

\subsubsection{Altitude Estimation}

Barometric altitude is available at \SI{50}{\hertz} and tracks fast vertical
motion, but drifts over the duration of a landing. A downward rangefinder is
accurate in the mean but becomes noisy within a rotor diameter of the ground and
reports at \SI{20}{\hertz}. The two are combined in a complementary filter that
propagates the barometric increment and corrects toward the rangefinder,
\begin{equation}
  \hat{h}_k = (1-\kappa)\,\hat{h}_k^{-} + \kappa\, h^{\mathrm{rf}}_k ,
  \qquad
  \kappa = \frac{\Delta t}{\tau + \Delta t},
  \label{eq:compfilter}
\end{equation}
where $\hat{h}_k^{-}$ is the barometric prediction, $h^{\mathrm{rf}}_k$ is the
rangefinder measurement, $\Delta t$ is the elapsed interval, and $\tau$ is the
blending time constant. When no rangefinder sample is available the estimate
propagates on the barometric prediction alone, so that dropouts degrade the
estimate gradually rather than interrupting it.

\subsubsection{Visual Alignment}

An AprilTag detector supplies the tag center in pixel coordinates
$\mathbf{p} = [u, v]^{\mathsf{T}}$ over a local publish--subscribe transport.
Detections arrive at \SI{8}{\hertz} while setpoints are streamed at
\SI{50}{\hertz}, so the controller does not consume detections synchronously.
Each arriving detection updates an exponential moving average,
\begin{equation}
  \bar{\mathbf{p}}_k = \alpha\, \mathbf{p}_k + (1-\alpha)\, \bar{\mathbf{p}}_{k-1},
  \label{eq:ema}
\end{equation}
and the smoothed estimate is held between detections, which both attenuates
detector jitter and supplies a value at every control cycle.

The image error is taken with respect to the principal point
$\mathbf{p}_c = [u_c, v_c]^{\mathsf{T}}$ and normalized by the half-width and
half-height of the image, making the error dimensionless and equal in scale on
both axes. Lateral velocity setpoints follow a proportional--derivative law in
which the derivative term acts on measured body velocity rather than on a
differentiated image error,
\begin{equation}
  v_x = -K_p \frac{v - v_c}{v_c} - K_d\, \tilde{v}_x ,
  \qquad
  v_y = -K_p \frac{u - u_c}{u_c} - K_d\, \tilde{v}_y ,
  \label{eq:pd}
\end{equation}
where $\tilde{v}_x$ and $\tilde{v}_y$ are the estimated body-frame velocities.
Differentiating the \SI{8}{\hertz} image signal would amplify detector noise at
the \SI{50}{\hertz} setpoint rate, whereas the onboard velocity estimate is
already available at full rate. Both components are saturated at
$v^{\max}_{xy}$ before being issued.

\subsubsection{Descent Sequencing}

Vertical motion is gated on lateral convergence: descent is commanded only when
the smoothed tag center lies within a radius $r_g$ of the principal point, and is
suspended whenever the error leaves that radius. Below a flare altitude
$h_f$ the descent rate is reduced from $v^{\max}_{z}$ to the touchdown rate
$v_{\mathrm{td}}$ for contact. The supervisory state machine sequences
\textsc{search}, \textsc{align}, \textsc{descend}, \textsc{flare}, and
\textsc{done} (Fig.~\ref{fig:statemachine}).

Tag detection is not guaranteed at every cycle. If no detection has been received
for longer than a timeout $T_{\mathrm{to}}$, the controller issues the last valid
setpoint scaled by one half, so that the vehicle coasts along its previous
command with decaying authority instead of either freezing at full command or
dropping to zero. If misses persist beyond $n_{\max}$ consecutive cycles the
supervisor reverts to \textsc{search}. Controller parameters are listed in
Table~\ref{tab:params}.

\subsection{Experimental Setup}

Experiments used a \SI{450}{\milli\meter} quadrotor running PX4~1.14 with a
downward-facing global-shutter camera at $640 \times 480$ pixels and 8 frames
per second. The landing platform was a treadmill-driven cart with a
\SI{300}{\milli\meter}-radius deck translating at \SI{0.3}{\meter\per\second}, and
carried a \SI{160}{\milli\meter} AprilTag of the 36h11 family. An OptiTrack
motion-capture system sampling at \SI{120}{\hertz} provided ground truth for
both vehicle and platform, from which the touchdown offset was computed.

Two controller configurations were compared. The proposed configuration
(\textit{vision}) uses the visually servoed law of \eqref{eq:pd} with the gated
descent described above. The baseline configuration (\textit{baro}) disables the
visual channel and commands an open-loop descent at $v^{\max}_{z}$, relying on
altitude estimation alone; it represents landing without platform feedback. Both
configurations share the same altitude estimator, setpoint rate, and airframe, so
the comparison isolates the visual channel.

Twelve landings were flown in each configuration ($N=12$ per mode), alternated
between configurations so that battery state and ambient wind were distributed
across both. The primary outcome was the touchdown offset, defined as the radial
distance in the deck plane between the vehicle contact point and the deck center
at the instant of contact. Secondary outcomes were the descent time, measured
from the start of the descent phase to contact, and the landing success rate,
defined as the proportion of trials with touchdown offset inside the
\SI{300}{\milli\meter} deck radius.

\subsection{Statistical Analysis}

Touchdown offset and descent time are reported as mean $\pm$ standard deviation
over the twelve trials of each configuration. Because trials were flown in
alternating pairs under matched conditions, the two configurations were compared
with a two-sided paired $t$-test on the touchdown offset, with significance at
$\alpha = 0.05$. The paired mean difference is reported with its
\SI{95}{\percent} confidence interval, since the interval conveys the magnitude
of the improvement in millimeters whereas the $p$-value alone does not. Success
rate is reported as a proportion of trials; no inferential test is applied to it,
as twelve trials per configuration do not support one. Trials were additionally
grouped by heading relative to the prevailing wind; that split is reported
descriptively only, because the resulting subgroups are too small to test.

\section{Results}

\subsection{Touchdown Accuracy}

The visually servoed configuration placed the vehicle within
$46 \pm \SI{18}{\milli\meter}$ of the deck center, against
$187 \pm \SI{92}{\milli\meter}$ for the altitude-only baseline
(Table~\ref{tab:results}). The paired mean difference was
\SI{141}{\milli\meter}, \SI{95}{\percent} CI $[98, 184]$~mm,
$p = 0.0002$. The confidence interval excludes zero over its full width and its
lower bound alone is a third of the deck radius.

The two configurations also differ in spread, not only in center: the baseline
standard deviation is five times that of the visual configuration
(\SI{92}{\milli\meter} against \SI{18}{\milli\meter}). Fig.~\ref{fig:offsets}
shows the per-trial offsets.

\subsection{Landing Success}

All twelve visually servoed landings finished inside the
\SI{300}{\milli\meter} deck radius. Three of twelve baseline landings did not,
that is, a success rate of \SI{100}{\percent} against \SI{75}{\percent}. The
baseline failures are consistent with its dispersion: a mean of
\SI{187}{\milli\meter} with a standard deviation of \SI{92}{\milli\meter} places
the deck edge slightly beyond one standard deviation of the baseline
distribution.

\subsection{Descent Time}

Visually servoed landings took $14.2 \pm \SI{1.9}{\second}$ from the start of
descent to contact, against $9.8 \pm \SI{0.7}{\second}$ for the baseline. The
visual configuration is therefore \SI{4.4}{\second} slower on average, and its
trial-to-trial variation is larger. This follows from the descent gate: vertical
motion is suspended whenever the lateral error leaves the gate radius, so the
duration of a landing depends on how often the platform motion pushes the
alignment out of tolerance, whereas the open-loop baseline descends at a fixed
rate irrespective of alignment.

\subsection{Sensitivity to Wind Heading}

Within the visually servoed trials, offsets were
$41 \pm \SI{15}{\milli\meter}$ downwind and $52 \pm \SI{20}{\milli\meter}$
upwind. The two groups overlap within one standard deviation, and both remain
well inside the deck radius.

\subsection{Behavior Under Detection Dropout}

One visually servoed trial contained two detection gaps exceeding the
\SI{0.4}{\second} timeout. In both, the decaying-command hold engaged and the
controller re-acquired the tag and completed the landing within the deck radius.
No trial in either configuration was aborted.

% ---------------------------------------------------------------------------
% Tables
% ---------------------------------------------------------------------------

\begin{table}[!t]
\caption{Controller Parameters}
\label{tab:params}
\centering
\small
\setlength{\tabcolsep}{4pt}
\begin{tabular}{llc}
\toprule
Parameter & Symbol & Value \\
\midrule
Setpoint rate                 & --                    & \SI{50}{\hertz} \\
Tag detection rate            & --                    & \SI{8}{\hertz} \\
Smoothing coefficient         & $\alpha$              & 0.35 \\
Lateral proportional gain     & $K_p$                 & 0.9 \\
Lateral derivative gain       & $K_d$                 & 0.15 \\
Lateral speed limit           & $v^{\max}_{xy}$       & \SI{1.2}{\meter\per\second} \\
Vertical speed limit          & $v^{\max}_{z}$        & \SI{0.5}{\meter\per\second} \\
Descent gate radius           & $r_g$                 & 40 px \\
Detection timeout             & $T_{\mathrm{to}}$     & \SI{0.4}{\second} \\
Consecutive-miss limit        & $n_{\max}$            & 3 \\
Flare altitude                & $h_f$                 & \SI{0.35}{\meter} \\
Touchdown descent rate        & $v_{\mathrm{td}}$     & \SI{0.25}{\meter\per\second} \\
Altitude blending constant    & $\tau$                & \SI{0.8}{\second} \\
\bottomrule
\end{tabular}
\end{table}

\begin{table}[!t]
\caption{Landing Performance Over Twelve Trials Per Configuration}
\label{tab:results}
\centering
\small
\setlength{\tabcolsep}{4pt}
\begin{tabular}{lccc}
\toprule
Metric & Vision & Baseline & $p$ \\
\midrule
Touchdown offset (mm)   & $46 \pm 18$    & $187 \pm 92$  & 0.0002 \\
Descent time (s)        & $14.2 \pm 1.9$ & $9.8 \pm 0.7$ & -- \\
Landings on deck        & 12/12          & 9/12          & -- \\
\bottomrule
\end{tabular}
\end{table}
